\documentclass[11pt]{article}
\usepackage[margin=1in]{geometry}
\usepackage[T1]{fontenc}
\usepackage[utf8]{inputenc}
\usepackage{amsmath,amssymb,amsthm}
\usepackage{enumitem}

\title{ICPC World Finals 2022\\U. Toy Train Tracks}
\author{}
\date{}

\begin{document}
\maketitle

\section*{Problem Summary}

We are given at most $s$ straight track pieces and at most $c$ curved pieces. A valid answer is a single
simple closed loop on the square grid. We must output any loop of maximum possible length.

\section*{Key Observations}

\begin{itemize}[leftmargin=*]
    \item In every closed orthogonal loop, the number of straight pieces is even and the number of curved
    pieces is even. So we may discard one unused piece of either type if its count is odd.
    \item If no straight pieces are used, the best possible lengths are:
    \begin{itemize}[leftmargin=*]
        \item $4$, using \texttt{LLLL};
        \item or any multiple of $4$ at least $12$, using a four-sided staircase loop.
    \end{itemize}
    There is no all-curves loop of length $8$.
    \item If at least two straight pieces are available, then every even number of curves at least $4$ is
    achievable together with every even number of straight pieces.
    \item The implementation constructs the loop as a sequence of \emph{absolute movement directions} on
    the grid and converts that cyclic direction sequence into the required alphabet
    \texttt{S/L/R} by comparing consecutive directions.
\end{itemize}

\section*{Algorithm}

\begin{enumerate}[leftmargin=*]
    \item Replace $(s,c)$ by the largest even counts not exceeding them.
    \item If fewer than two straight pieces remain:
    \begin{itemize}[leftmargin=*]
        \item use $4$ curves if fewer than $12$ curves remain;
        \item otherwise use the largest multiple of $4$ not exceeding $c$.
    \end{itemize}
    \item If at least two straight pieces remain and $c \equiv 0 \pmod 4$, build the direction sequence
    \[
    N^a\ W\ (SW)^q\ S^a\ E\ (NE)^q,
    \]
    where $a=s/2+1$ and $q=(c-4)/4$.
    \item If at least two straight pieces remain and $c \equiv 2 \pmod 4$, build
    \[
    N^2\ (WN)^p\ W^b\ (SE)^{p+1}\ S\ E^{b-1},
    \]
    where $b=s/2+1$ and $p=(c-6)/4$.
    \item Convert the cyclic direction sequence into \texttt{S/L/R} and output it.
\end{enumerate}

\section*{Correctness Proof}

We prove that the algorithm returns the correct answer.

\paragraph{Lemma 1.}
Any closed loop uses an even number of straight pieces and an even number of curved pieces.

\paragraph{Proof.}
The total horizontal displacement and the total vertical displacement of a closed loop are both zero, so
the number of steps in each direction must balance. This forces the number of straight pieces to be even.
Also, the total turning angle of a simple closed loop is $\pm 360^\circ$, so the difference between the
numbers of left and right turns is $\pm 4$. Therefore their sum, which is the total number of curved
pieces, is even. \qed

\paragraph{Lemma 2.}
The direction families used by the algorithm are simple closed loops and use exactly the claimed numbers
of straight and curved pieces.

\paragraph{Proof.}
Each family is a concatenation of monotone staircase sides. Their horizontal and vertical displacements
cancel by construction, so the walk closes. Because the sides are monotone and nested outward, the walk
never revisits a cell before the end, so it is simple.

Counting consecutive equal directions gives the number of straight pieces, and counting direction changes
gives the number of curved pieces. For the two families these counts are respectively
\[
2a-2,\ 4+4q
\quad\text{and}\quad
2b-2,\ 6+4p,
\]
which are exactly the requested values after substituting $a=s/2+1$, $b=s/2+1$, $q=(c-4)/4$, and
$p=(c-6)/4$. \qed

\paragraph{Lemma 3.}
For the all-curves case, the algorithm uses the largest feasible number of curves.

\paragraph{Proof.}
By Lemma 1, only even numbers of curves are possible. The staircase all-curves family realizes every
multiple of $4$ at least $12$, and \texttt{LLLL} realizes $4$. The only even value below $12$ not covered
by these constructions is $8$, which is impossible. Hence the algorithm picks the largest feasible number
of curves when no straight pieces can be used. \qed

\paragraph{Theorem.}
The algorithm outputs a loop of maximum possible length.

\paragraph{Proof.}
By Lemma 1, any optimal solution can use only even numbers of straight and curved pieces, so rounding
the available counts down to even values loses nothing.

If fewer than two straight pieces remain, Lemma 3 shows that the algorithm uses the maximum possible
number of curves. Otherwise, Lemma 2 shows that the algorithm can realize every even curve count at
least $4$ together with every even straight count, so it uses all available pieces after the parity adjustment.
No longer loop can exist. \qed

\section*{Complexity Analysis}

The algorithm constructs the answer string directly. The running time is linear in the output length, and
the memory usage is also linear in the output length.

\section*{Implementation Notes}

\begin{itemize}[leftmargin=*]
    \item The output does not need to match the sample exactly; any maximum-length simple loop is
    accepted.
\end{itemize}

\end{document}
